%!TEX program = xelatex
%!TEX encoding = UTF-8
%Defini\c{c}\~oes de Ano, Semestre, n\'umero e data da avalia\c{c}\~ao
\def\numeroteste{1}
\def\modulo{1}
\def\dataavaliacao{06/09/2023}

\documentclass[12pt]{exam}

\usepackage[singlelinecheck=false]{caption}
\usepackage{textpos}
\usepackage{amssymb}
\usepackage{amsmath,amsfonts,amsthm,amstext}
\usepackage[justification=centering]{caption}
\usepackage[brazil]{babel}
% \usepackage[latin1]{inputenc}
\usepackage{graphicx}
\graphicspath{{../../figuras_geogebra/}{../../Pictures/}}
\usepackage{enumitem}
\usepackage{multicol}
\usepackage{answers}
\usepackage{qrcode}

\usepackage{pifont}
\newcommand{\cmark}{\ding{228}}%

% requires version 0.3 of the package
\usepackage[customcolors,beamer]{hf-tikz}

\tikzset{style green/.style={
        set fill color=green!50!lime!60,
        set border color=white,
    },
    style cyan/.style={
        set fill color=cyan!90!blue!60,
        set border color=white,
    },
    style orange/.style={
        set fill color=orange!80!red!60,
        set border color=white,
    },
    hor/.style={
        above left offset={-0.15,0.31},
        below right offset={0.15,-0.125},
        #1
    },
    ver/.style={
        above left offset={-0.1,0.3},
        below right offset={0.15,-0.15},
        #1
    }
}

\usepackage{tikz,ifthen,tikz-3dplot}
\usepackage{siunitx}
\usetikzlibrary{circuits.ee.IEC}
\usetikzlibrary{lindenmayersystems}
\usetikzlibrary[shadings]
\usepackage{pgfplots}
\pgfplotsset{compat=1.15}
\usepackage{mathrsfs}
\usetikzlibrary{arrows}
\usepackage{tcolorbox}
\newcounter{saveenumi}
\newcommand{\seti}{\setcounter{saveenumi}{\value{enumi}}}
\newcommand{\conti}{\setcounter{enumi}{\value{saveenumi}}}

\resetcounteronoverlays{saveenumi}


\renewcommand{\qedsymbol}{$\blacksquare$}

%\setbeamertemplate{itemize items}[ball]

%\setbeamertemplate{itemize items}[default]
%-------------amatrix
% Augmented matrix.  Usage (note the argument does not count the aug col):
% \begin{amatrix}{2}
    %   1  2  3 \\  4  5  6
    % \end{amatrix}
\newenvironment{amatrix}[1]{%
    \left[\begin{array}{@{}*{#1}{r}|r@{}}
    }{%
    \end{array}\right]
}

\def\ano{2023}
\def\semestre{2}
\def\turma{11}
\newcommand{\nsub}{\varsubsetneq}
\newcommand{\vaz}{\emptyset}
\newcommand{\dom}{{\rm dom\,}}
\newcommand{\im}{{\rm Im\,}}
\newcommand{\sub}{\subseteq}
\newcommand{\n}{\mathbb{N}}
\newcommand{\z}{\mathbb{Z}}
\newcommand{\rac}{\mathbb{Q}}
\newcommand{\real}{\mathbb{R}}
\newcommand{\complex}{\mathbb{C}}
\newcommand{\cp}[1]{\mathbb{#1}}
\newcommand{\ch}{\mbox{\textrm{car\,}}\nobreak}
\newcommand{\p}[1]{\mbox{\textrm{posto(#1)}}}
\newcommand{\vesp}[1]{\vspace{ #1  cm}}
\newcommand{\compcent}[1]{\vcenter{\hbox{$#1\circ$}}}
\newcommand{\comp}{\mathbin{\mathchoice
{\compcent\scriptstyle}{\compcent\scriptstyle}
{\compcent\scriptscriptstyle}{\compcent\scriptscriptstyle}}}
\newcommand\T{\rule{0pt}{2.6ex}}
\newtheorem{definicao}{Defini\c{c}\~ao}[section]
\newtheorem{definicoes}{Defini\c{c}\~oes}[section]
\newtheorem{exemplo}{Exemplo}[section]
\newtheorem{exemplos}{Exemplos}[section]
\newtheorem{exercicio}{Exerc{\'\i}cio}
\newtheorem{observacao}{Observa{\c c}{\~a}o:}[section]
\newtheorem{observacoes}{Observa{\c c}{\~o}es:}[section]
\newtheorem*{solucao}{Solu{\c c}{\~a}o:}
\newtheorem{proposicao}{Proposi\c{c}\~ao}
\newtheorem{lema}{Lema}
\newtheorem{teorema}{Teorema}
\newtheorem{corolario}{Corol\'ario}
\newenvironment{prova}[1][Prova]{\noindent\textbf{#1:} }{\qedsymbol}%{\ \rule{0.5em}{0.5em}}

\begin{document}

    \begin{figure}[h]
        \begin{minipage}[c]{1.7cm}
            \includegraphics[width=1.7cm]{unb.pdf}
        \end{minipage}
        \hspace{0pt}
        \begin{minipage}[c]{4in}
            {Universidade de Bras{\'\i}lia} \\
            {Departamento de Matem{\'a}tica}
        \end{minipage}
    \end{figure}
    \vesp{-1}
    \hrule
    \begin{center}
        {\Large\bf \disciplina\ - Turma \turma}  \\
        \vesp{0.2} {\large\bf Teste \numeroteste\ - Módulo\ \modulo\ -\ \dataavaliacao}
    \end{center}

    \noindent{\bf \underline{Nome: \hspace{9.8cm} Mat.:\qquad \
            \hspace{2.1cm}}}\\
    \vspace*{.01cm}

    \noindent{\bf \underline{Nome: \hspace{9.8cm} Mat.:\qquad \
            \hspace{2.1cm}}}
    \vspace{.4cm}

    \questao{} Considere o seguinte sistema linear:
    \[
    \begin{cases}
        -2x_1 + 5x_2 - 10x_3 = 4\\
        x_1 - 2x_2 + 3x_3 = -1\\
        7x_1 - 17x_2 + 34x_3 = -16
    \end{cases}
    \]
    \begin{enumerate}[label={\alph*})]
        \item Escreva a matriz aumentada desse sistema.

        \item Encontre a solução desse sistema aplicando os métodos vistos em sala.
    \end{enumerate}
    \solucao
    \begin{enumerate}[label={\alph*})]
        \item A matriz ampliada é:
        \[
            \begin{amatrix}{3}
                -2 & 5 & -10 & 4\\
                1 & -2 & 3 & -1\\
                7 & -17 & 34 & -16
            \end{amatrix}
        \]

        \item Aplicando as operações elementares até transformar a matriz ampliada numa matriz escalonada reduzida por linhas:
        \begin{align*}
            &\begin{amatrix}{3}
                -2 & 5 & -10 & 4\\
                1 & -2 & 3 & -1\\
                7 & -17 & 34 & -16
            \end{amatrix}
            \begin{array}{l}
                L_1 \leftrightarrow L_2\\\phantom{x}\\\phantom{x}
            \end{array}\sim
            &\begin{amatrix}{3}
                1 & -2 & 3 & -1\\
                -2 & 5 & -10 & 4\\
                7 & -17 & 34 & -16
            \end{amatrix}
            \begin{array}{l}
                \phantom{x}\\L_2 \to L_2 + 2L_1\\L_3 \to L_3 - 7L_1
            \end{array}\sim\\
            &\begin{amatrix}{3}
                1 & -2 & 3 & -1\\
                0 & 1 & -4 & 1\\
                0 & -3 & 13 & -9
            \end{amatrix}
            \begin{array}{l}
                \phantom{x}\\\phantom{x}\\L_3 \to L_3 + 3L_2
            \end{array}\sim
            &\begin{amatrix}{3}
                1 & -2 & 3 & -1\\
                0 & 1 & -4 & 1\\
                0 & 0 & 1 & -3
            \end{amatrix}
            \begin{array}{l}
                \phantom{x}\\L_2 \to L_2 + 4L_3\\\phantom{x}
            \end{array}\sim\\
            &\begin{amatrix}{3}
                1 & -2 & 3 & -1\\
                0 & 1 & 0 & -10\\
                0 & 0 & 1 & -3
            \end{amatrix}
            \begin{array}{l}
                L_1 \to L_1 - 3L_3\\\phantom{x}\\\phantom{x}
            \end{array}\sim
            &\begin{amatrix}{3}
                1 & -2 & 0 & 8\\
                0 & 1 & 0 & -10\\
                0 & 0 & 1 & -3
            \end{amatrix}
            \begin{array}{l}
                L_1 \to L_1 + 2L_2\\\phantom{x}\\\phantom{x}
            \end{array}\sim\\
            &\begin{amatrix}{3}
                1 & 0 & 0 & -12\\
                0 & 1 & 0 & -10\\
                0 & 0 & 1 & -3
            \end{amatrix}
        \end{align*}
        Portanto a solução do sistema é $x_1 = -12$, $x_2 = -10$ e $x_3 = -3$.
    \end{enumerate}
\end{document}
